%------------------------------------------------------------------------------%
\documentclass[a4paper,12pt,fleqn]{article}

\usepackage{gaal}

\ShowAnswers

%------------------------------------------------------------------------------%
\begin{document}

\makeHeader{GAAL}{Prova 4 -- A}{2026-2}

\bigskip
Dada a matriz
\(
  B = \begin{pmatrix}
    4 & 0 & 0 \\
    0 & 3 & 1 \\
    0 & 1 & 3
  \end{pmatrix}
\)

% 01
%------------------------------------------------------------------------------%
\exercise{25\%}
Determine os autovalores de $B$

\begin{answer}
  Polinômio característico
  \begin{align*}
    p(\lambda)
    & = \det(B-\lambda I) \\
    & = (-1)^{1+1}(4-\lambda)
        \begin{vmatrix}
          3-\lambda & 1         \\
          1         & 3-\lambda
        \end{vmatrix}\\
    & = (4-\lambda) \left[(3-\lambda)^2-1\right] \\
    & = (4-\lambda)(9-6\lambda +\lambda^2 - 1) \\
    & = (4-\lambda)(\lambda^2 - 6\lambda + 8)
  \end{align*}
  Temos \(\lambda=4\) ou \(\lambda^2 - 6\lambda + 8=0\)
  \[
    \Delta
    = b^2 - 4ac
    = (-6)^2 - 4 \times 1 \times 8
    = 36 - 32
    = 4
  \]
  \[
    \sqrt{\Delta} = \sqrt{4} = 2
  \]
  \[
    \lambda
    = \frac{-b - \sqrt{\Delta}}{2a}
    = \frac{6+2}{2}
    = 4
  \]
  \[
    \lambda
    = \frac{-b - \sqrt{\Delta}}{2a}
    = \frac{6-2}{2}
    = 2
  \]
  O autovalor $4$ possui multiplicidade algébrica dois, isto é,
  ele aparece duas vezes como solução da equação característica
  \[
    \lambda_1 = 2 \qquad
    \lambda_2 = 4 \qquad
    \lambda_3 = 4
  \]
\end{answer}

% 02
%------------------------------------------------------------------------------%
\exercise{25\%}
Determine os autovetores e autoespaços de $B$

\begin{answer}
\begin{multicols}{2}
\textbf{Autoespaço associado à $\lambda_1=2$}
  \begin{align*}
    (B-2I)\vec{v} & = \vec{0} \\
    \begin{pmatrix}
      2 & 0 & 0 \\
      0 & 1 & 1 \\
      0 & 1 & 1
    \end{pmatrix}
    \vetortri{x}{y}{z} & = \vetortri{0}{0}{0}
  \end{align*}
  O sistema é
  \begin{align*}
    2x  & = 0 \\
    y+z & = 0
  \end{align*}
  com solução
  \(
    x=0
  \) e
  \(
    z=-y
  \)

  tomando $y=1$ temos o autovetor
  \[
    \vec{v}_1 = \vetortri{0}{1}{-1}
  \]
  e o autoespaço
  \[
    E(2) = \ger \left\{ \vetortri{0}{1}{-1} \right\}
  \]

\bigskip
\textbf{Autoespaço associado à $\lambda_2=\lambda_3=4$}
  \begin{align*}
    (B-4I)\vec{v} & = \vec{0} \\
    \begin{pmatrix}
      0 & 0  & 0  \\
      0 & -1 & 1  \\
      0 & 1  & -1
    \end{pmatrix}
    \vetortri{x}{y}{z} & = \vetortri{0}{0}{0}
  \end{align*}
  O valor de $x$ é livre e temos
  \(
    -y + z = 0
  \) que corresponde a
  \(
    z = y
  \).
  A solução do sistema é
  \[
    \vetortri{x}{y}{z}
    = x \vetortri{1}{0}{0}
    + y \vetortri{0}{1}{1}
  \]
  Os dois autovetores são
  \[
    \vec{v}_2 = \vetortri{1}{0}{0} \qquad
    \vec{v}_3 = \vetortri{0}{1}{1}
  \]
  O autoespaço é
  \[
    E(4) = \ger \left\{ \vetortri{1}{0}{0}, \vetortri{0}{1}{1} \right\}
  \]
\end{multicols}
\end{answer}

% 03
%------------------------------------------------------------------------------%
\exercise{25\%}
Diagonalize ortogonalmente a matriz $B$

\begin{answer}
  A matriz $B$ é simétrica, portanto poder ser diagonalizada ortogonalmente

  Precisamos normalizar os autovetores de $B$

  Normas
  \begin{align*}
    \norm{\vec{v}_1}
    & = \norm{\vetortri{0}{1}{-1}}
    = \sqrt{0^2 + 1^2 + (-1)^2}
    = \sqrt{2}
    \\[5mm]
    \norm{\vec{v}_2}
    & = \norm{\vetortri{1}{0}{0}}
    = \sqrt{1^2 + 0^2 + 0^2}
    = 1
    \\[5mm]
    \norm{\vec{v}_3}
    & = \norm{\vetortri{0}{1}{1}}
    = \sqrt{1^2 + 1^2 + 0^2}
    = \sqrt{2}
  \end{align*}

  Autovetores normalizados
  \begin{align*}
    \vec{u}_1
    & = \frac{\vec{v}_1}{\norm{\vec{v}_1}}
    = \frac{1}{\sqrt{2}} \vetortri{0}{1}{-1}
    = \Vetortri{0}{\sfrac{1}{\sqrt{2}}}{\sfrac{-1}{\sqrt{2}}}
    \\[5mm]
    \vec{u}_2
    & = \frac{\vec{v}_2}{\norm{\vec{v}_2}}
    = \frac{1}{1} \vetortri{1}{0}{0}
    = \vetortri{1}{0}{0}
    \\[5mm]
    \vec{u}_3
    & = \frac{\vec{v}_3}{\norm{\vec{v}_3}}
    = \frac{1}{\sqrt{2}} \vetortri{0}{1}{1}
    = \Vetortri{0}{\sfrac{1}{\sqrt{2}}}{\sfrac{1}{\sqrt{2}}}
  \end{align*}

  Matriz ortogonal
  \[
    P
    = \begin{bmatrix}
      \vec{u}_1 & \vec{u}_2 & \vec{u}_3
    \end{bmatrix}
    = \begin{bmatrix}
      0                    & 1 & 0                   \\[3mm]
      \dfrac{1}{\sqrt{2}}  & 0 & \dfrac{1}{\sqrt{2}} \\[3mm]
      \dfrac{-1}{\sqrt{2}} & 0 & \dfrac{1}{\sqrt{2}}
    \end{bmatrix}
    = \frac{1}{\sqrt{2}} \begin{bmatrix}[c]
      0  & \sqrt{2} & 0 \\
      1  & 0        & 1 \\
      -1 & 0        & 1
    \end{bmatrix}
  \]
  Matriz diagonal
  \[
    D
    = \diag( \lambda_1, \lambda_2, \lambda_3 )
    = \begin{bmatrix}
      2 & 0 & 0 \\
      0 & 4 & 0 \\
      0 & 0 & 4
    \end{bmatrix}
  \]
\end{answer}

% 04
%------------------------------------------------------------------------------%
\exercise{25\%}
Mostre que as matrizes produzidas no exercício anterior realmente diagonalizam $B$

\begin{answer}
  Queremos mostrar que \(B = PDP^t\)
  \begin{align*}
    PDP^t
    & = \frac{1}{\sqrt{2}} \begin{bmatrix}[c]
      0  & \sqrt{2} & 0 \\
      1  & 0        & 1 \\
      -1 & 0        & 1
    \end{bmatrix}
    \begin{bmatrix}
      2 & 0 & 0 \\
      0 & 4 & 0 \\
      0 & 0 & 4
    \end{bmatrix} P^t
    \\
    & = \frac{1}{\sqrt{2}} \begin{bmatrix}[c]
       0  & 4\sqrt{2} & 0 \\
       2  & 0         & 4 \\
      -2  & 0         & 4
    \end{bmatrix} P^t
    \\
    & = \frac{1}{\sqrt{2}} \begin{bmatrix}[c]
       0  & 4\sqrt{2} & 0 \\
       2  & 0         & 4 \\
      -2  & 0         & 4
    \end{bmatrix}
    \frac{1}{\sqrt{2}} \begin{bmatrix}[c]
      0        & 1 & -1  \\
      \sqrt{2} & 0 &  0  \\
      0        & 1 &  1
    \end{bmatrix}
    \\
    & = \frac{1}{2} \begin{bmatrix}[c]
      0 \times 0  + 4\sqrt{2} \sqrt{2}  + 4 \times 0 &
      0 \times 1  + 4\sqrt{2} \times 0  + 0 \times 1 &
      0 (-1)      + 4\sqrt{2} \times 0  + 4 \times 0
      \\
      2 \times 0 + 0 \sqrt{2} + 4 \times 0 &
      2 \times 1 + 0 \times 0 + 4 \times 1 &
      2 (-1)     + 0 \times 0 + 4 \times 1
      \\
      (-2) 0   + 0 \sqrt{2} + 4 \times 0 &
      (-2) 1   + 0 \times 0 + 4 \times 1 &
      (-2)(-1) + 0 \times 0 + 4 \times 1
    \end{bmatrix}
    \\
    & = \frac{1}{2} \begin{bmatrix}[c]
      8 & 0 & 0 \\
      0 & 6 & 2 \\
      0 & 2 & 6
    \end{bmatrix}
    \\
    & = \begin{bmatrix}[c]
      4 & 0 & 0 \\
      0 & 3 & 1 \\
      0 & 1 & 3
    \end{bmatrix}
  \end{align*}
\end{answer}

%------------------------------------------------------------------------------%
\end{document}
%------------------------------------------------------------------------------%

